% Part: incompleteness
% Chapter: representability-in-q
% Section: c

\documentclass[../../../include/open-logic-section]{subfiles}

\begin{document}

\olfileid{inc}{req}{cee}
\olsection{函数集 $C$}

令 $C$ 为包含下列函数的最小函数集合：

\begin{enumerate}
\item $0$，
\item 后继函数，
\item 加法，
\item 乘法，
\item 投影函数，以及
\item 等词的特征函数 $\Char{=}$；
\end{enumerate}

并且在下列运算下封闭：

\begin{enumerate}
\item 复合运算，以及
\item 对正则函数应用的无界搜索运算。
\end{enumerate}

注意，这最后一条限制仅仅意味着：只有当结果为全函数时，才能使用 $\mu$ 运算。
将此定义与\emph{一般递归函数}的定义相比较：这里我们加入了加法、乘法和 $\Char{=}$，但舍弃了原始递归。

显然，$C$ 中的所有函数都是递归的，因为加法、乘法和 $\Char{=}$ 也是递归的。
我们将证明其逆命题亦成立；这相当于说，利用 $C$ 中的其他组成部分可以执行原始递归。

\end{document}